\documentclass{article}%
\usepackage[T1]{fontenc}%
\usepackage[utf8]{inputenc}%
\usepackage{lmodern}%
\usepackage{textcomp}%
\usepackage{lastpage}%
\usepackage{geometry}%
\geometry{margin=1in,headheight=14pt,headsep=25pt}%
\usepackage{hyperref}%
\usepackage{enumitem}%
\usepackage{fancyhdr}%
\usepackage{xcolor}%
\usepackage{url}%
\usepackage{breakurl}%
\usepackage{amsmath}%
\usepackage{amssymb}%
\usepackage{mathtools}%
\usepackage{amsthm}%
\usepackage{thmtools}%
\usepackage{float}%
\usepackage{subcaption}%
\usepackage[font=small,labelfont=bf]{caption}%
%

            \hypersetup{
                colorlinks=true,
                linkcolor=blue,
                filecolor=magenta,
                urlcolor=blue,
                breaklinks=true,
                bookmarks=true,
                pdfborder={0 0 0}
            }
            \urlstyle{same}
            \Urlmuskip=0mu plus 1mu  % URL breaking in nicer spots

            % Math configuration
            \everymath{\displaystyle}
            \setlength{\jot}{10pt}  % Increase spacing between equations

            \pagestyle{fancy}
            \fancyhf{}
            \rhead{Generated on \today}
            \cfoot{\thepage}

            % Better Unicode support
            \usepackage[utf8]{inputenc}
            \usepackage{newunicodechar}
            \setlength{\itemsep}{1em} % Adds spacing between items

            % Configure itemize settings globally
            \setlist[itemize]{nosep}
            \setlist[itemize]{leftmargin=*}

            % Remove section numbering
            \renewcommand{\thesection}{}
            \renewcommand{\thesubsection}{}

            % Adjust caption spacing around figures
            \setlength{\abovecaptionskip}{0pt}
            \setlength{\belowcaptionskip}{5pt}
        %
\lhead{Lecture: 04L1{-}Linear Algebra Primer}%
\title{04L1{-}Linear Algebra Primer}%
\author{Generated by Scribe.AI}%
\date{\today}%
%
\begin{document}%
\normalsize%
\maketitle%
\section{Slide 1}%
\label{sec:Slide1}%
\subsection{Content}%

\textbf{CS 24300: AI Basics}

Tony Bergstrom

Slides adapted from Yexiang Xue and Tony Bergstrom's original set
%
\subsection{Description}%
This slide is a title slide for a lecture on Artificial Intelligence Basics, likely part of a course called CS 24300.  It displays the course title, the instructor's name (Tony Bergstrom), and a credit line indicating that the slides are adapted from a previous set created by Yexiang Xue and Tony Bergstrom.  The background is a dark gray, and the text is light gray, making it easily readable.  The context is a standard introductory slide for a lecture, providing basic information about the course and its materials.%
\subsection{Figures}%
\begin{figure}[H]%
\centering%
\end{figure}

%
\newpage%
\section{Slide 2}%
\label{sec:Slide2}%
\subsection{Content}%

\textbf{This is not a Linear Algebra Class}

But we'll briefly touch on some things you should know upfront
%
\subsection{Description}%
This slide is a title slide, likely the first slide of a lecture. It states that the course is not a linear algebra course, but that some prerequisite linear algebra concepts will be briefly reviewed.  The text is simple and straightforward, clearly indicating the course's focus and the nature of the introductory material.%
\subsection{Figures}%
\begin{figure}[H]%
\centering%
\end{figure}

%
\newpage%
\section{Slide 3}%
\label{sec:Slide3}%
\subsection{Content}%

\textbf{We're not teaching Linear Algebra}

Linear Algebra is a necessary topic to understand how much of this computation works.
You should and will take a Linear Algebra course if you haven't already, but...

... many of you haven't yet.

So I will go over some basics to help with some context.

It's not enough - so look forward to more depth when you take Linear Algebra.
%
\subsection{Description}%
This slide, likely from a lecture on Artificial Intelligence Basics, acknowledges that Linear Algebra is a prerequisite for a deeper understanding of the computational aspects of the subject.  It explicitly states that the current lecture will not cover Linear Algebra in detail, but will provide some basic background to aid in understanding the context of the AI concepts being presented.  The instructor emphasizes that a dedicated Linear Algebra course is necessary for a thorough understanding of the subject matter.%
\subsection{Figures}%
\begin{figure}[H]%
\centering%
\end{figure}

%
\newpage%
\section{Slide 4}%
\label{sec:Slide4}%
\subsection{Content}%

\textbf{A Matrix}

A matrix is a rectangular organization of numbers, symbols, or expressions of some dimension $n \times m$ such that it has $n$ columns and $m$ rows

$A = \begin{pmatrix}
a_{11} & a_{12} & \cdots & a_{1n} \\
a_{21} & a_{22} & \cdots & a_{2n} \\
\vdots & \vdots & \ddots & \vdots \\
a_{m1} & a_{m2} & \cdots & a_{mn}
\end{pmatrix}$
%
\subsection{Description}%
This slide defines a matrix.  It states that a matrix is a rectangular arrangement of numbers, symbols, or expressions.  The matrix has dimensions $n \times m$, meaning it has $n$ columns and $m$ rows.  The slide then presents the general form of a matrix, using subscripts to denote the elements.  The elements are labeled as $a_{ij}$, where $i$ represents the row number and $j$ represents the column number.  The context of the course (likely Linear Algebra or a course that uses matrices) is crucial for understanding the importance of matrices in mathematical modeling and computation.%
\subsection{Figures}%
\begin{figure}[H]%
\centering%
\end{figure}

%
\newpage%
\section{Slide 5}%
\label{sec:Slide5}%
\subsection{Content}%

\textbf{So what does a matrix represent?}

A Linear Transform

$T: \mathbb{R}^2 \rightarrow \mathbb{R}^2$

$T(x) = \begin{pmatrix} 0 & -1 \\ 1 & 0 \end{pmatrix}$

As a transform, we are mapping some space to other points to other points in space.

https://textbooks.math.gatech.edu/ila/linear-transformations.html
%
\subsection{Description}%
This slide, likely from a lecture on Artificial Intelligence Basics, introduces the concept of a matrix representing a linear transformation.  It shows a linear transformation $T$ from $\mathbb{R}^2$ to $\mathbb{R}^2$ represented by a $2 \times 2$ matrix. The matrix is $\begin{pmatrix} 0 & -1 \\ 1 & 0 \end{pmatrix}$. The slide emphasizes that a matrix can be used to represent a linear transformation, which maps points in one space to points in another space. The URL at the bottom of the slide likely points to a resource explaining linear transformations and matrices in more detail.  The context of the course is crucial for understanding how matrices are used to represent transformations in AI algorithms, such as image processing or data manipulation.%
\subsection{Figures}%
\begin{figure}[H]%
\centering%
\end{figure}

%
\newpage%
\section{Slide 6}%
\label{sec:Slide6}%
\subsection{Content}%

\textbf{But it is a Linear Transform...}

As its a linear transform, this means $T: \mathbb{R}^n \rightarrow \mathbb{R}^m$ satisfies

$T(u + v) = T(u) + T(v)$

$T(cu) = cT(u)$

For all vectors $u, v$ and scalars $c$

If we represent $T(x) = Ax$, using matrix multiplication to apply this transform, we have:

$T(u + v) = A(u + v) = Au + Av = T(u) + T(v)$

$T(cu) = A(cu) = cAu = cT(u)$

https://textbooks.math.gatech.edu/ila/linear-transformations.html
%
\subsection{Description}%
This slide defines a linear transformation, $T$, from $\mathbb{R}^n$ to $\mathbb{R}^m$.  It states that a linear transformation must satisfy two properties:  $T(u + v) = T(u) + T(v)$ and $T(cu) = cT(u)$ for all vectors $u, v$ in $\mathbb{R}^n$ and scalar $c$.  The slide then shows how these properties are reflected when the transformation is represented by matrix multiplication.  If $T(x) = Ax$, then $T(u + v) = A(u + v) = Au + Av = T(u) + T(v)$ and $T(cu) = A(cu) = cAu = cT(u)$.  The URL at the bottom of the slide likely points to a resource explaining linear transformations and matrices in more detail.  The context of the course (Artificial Intelligence Basics) suggests that this slide is part of a section introducing the mathematical underpinnings of AI, specifically how linear transformations are represented and used in various AI algorithms.%
\subsection{Figures}%
\begin{figure}[H]%
\centering%
\end{figure}

%
\newpage%
\section{Slide 7}%
\label{sec:Slide7}%
\subsection{Content}%

\textbf{Examples of Linear Transforms}

Rotation \\
Scaling \\
Projections \\
Essentially it's the mapping of points, lines, planes, spaces to other points, lines, planes, spaces.
%
\subsection{Description}%
This slide, likely from a lecture on Artificial Intelligence Basics, lists examples of linear transformations.  It explicitly names "Rotation," "Scaling," and "Projections" as examples.  The slide then provides a concise definition of a linear transformation as the mapping of points, lines, planes, and spaces to other points, lines, planes, and spaces.  The context of the course suggests that this slide is part of a section introducing the mathematical underpinnings of AI, specifically how linear transformations are used in various AI algorithms, such as image processing or data manipulation.  The examples given are fundamental geometric transformations that can be represented by matrices.%
\subsection{Figures}%
\begin{figure}[H]%
\centering%
\end{figure}

%
\newpage%
\section{Slide 8}%
\label{sec:Slide8}%
\subsection{Content}%
%
\subsection{Description}%
This slide in a course on Artificial Intelligence Basics discusses further properties of linear transformations. It states that a linear transformation, represented by a matrix, maps the zero vector to the zero vector ($T(\vec{0}) = \vec{0}$).  Crucially, it highlights the linearity property:  a linear transformation applied to a linear combination of vectors is the same as the linear combination of the transformations applied to each vector individually.  This is formally expressed as $T(c_1v_1 + c_2v_2 + \dots + c_kv_k) = c_1T(v_1) + c_2T(v_2) + \dots + c_kT(v_k)$.  The slide also introduces the concept of combining linear transformations, $T$ and $U$, through addition, $(T+U)(x) = T(x) + U(x)$.  Finally, it emphasizes that a linear transformation can reduce the dimensionality of a space but cannot increase it. The URL at the bottom likely points to a resource providing further details on linear transformations.%
\subsection{Figures}%
\begin{figure}[H]%
\centering%
\end{figure}

%
\newpage%
\section{Slide 9}%
\label{sec:Slide9}%
\subsection{Content}%

\textbf{Adding, subtracting, and scalar multiplication}

Addition and Subtraction are element-wise operations
the matrices should be the same dimensions

$A + B = C$ where $c_{ij} = a_{ij} + b_{ij}$

$\begin{pmatrix}
1 & 2 \\
3 & 4
\end{pmatrix} + \begin{pmatrix}
5 & 6 \\
7 & 8
\end{pmatrix} = \begin{pmatrix}
1 + 5 & 2 + 6 \\
3 + 7 & 4 + 8
\end{pmatrix} = \begin{pmatrix}
6 & 8 \\
10 & 12
\end{pmatrix}$

Scalar multiplication multiplies each element

$k \cdot A = B$ where $b_{ij} = k \cdot a_{ij}$

$2 \cdot \begin{pmatrix}
1 & 2 \\
3 & 4
\end{pmatrix} = \begin{pmatrix}
2 \cdot 1 & 2 \cdot 2 \\
2 \cdot 3 & 2 \cdot 4
\end{pmatrix} = \begin{pmatrix}
2 & 4 \\
6 & 8
\end{pmatrix}$
%
\subsection{Description}%
This slide from a course on Artificial Intelligence Basics covers matrix operations: addition, subtraction, and scalar multiplication.  It defines these operations for matrices.  Addition and subtraction are element-wise, meaning corresponding elements of the matrices are added or subtracted.  The matrices involved in these operations must have the same dimensions.  The slide provides an example of matrix addition, demonstrating how corresponding elements are added to produce the resulting matrix.  Similarly, scalar multiplication is shown, where each element of a matrix is multiplied by a scalar value.  The example shows how a matrix is multiplied by a scalar 2.  The context of the course suggests that this slide is part of a section introducing the mathematical tools used in AI, specifically matrix operations, which are fundamental to many AI algorithms.%
\subsection{Figures}%
\begin{figure}[H]%
\centering%
\end{figure}

%
\newpage%
\section{Slide 10}%
\label{sec:Slide10}%
\subsection{Content}%

\textbf{Inverses}

Some Matrices have inverses. All of those matrices are square.
Lots of Handwaving on details

$A = \begin{pmatrix} 1 & 2 \\ 3 & 4 \end{pmatrix}$
$A^{-1} = \begin{pmatrix} -2 & 1 \\ \frac{3}{2} & -\frac{1}{2} \end{pmatrix}$
$AA^{-1} = I = \begin{pmatrix} 1 & 0 \\ 0 & 1 \end{pmatrix}$

Note, because they are inverses of each other, $AA^{-1} = A^{-1}A = I$
%
\subsection{Description}%
This slide from a course on Artificial Intelligence Basics discusses matrix inverses. It shows an example of a 2x2 matrix $A$ and its inverse $A^{-1}$.  The slide states that some matrices have inverses, and all such matrices are square.  The example shows the matrix $A$ as $\begin{pmatrix} 1 & 2 \\ 3 & 4 \end{pmatrix}$ and its inverse $A^{-1}$ as $\begin{pmatrix} -2 & 1 \\ \frac{3}{2} & -\frac{1}{2} \end{pmatrix}$.  The slide then demonstrates that the product of the matrix and its inverse, $AA^{-1}$, is the identity matrix $I$, which is $\begin{pmatrix} 1 & 0 \\ 0 & 1 \end{pmatrix}$.  The final line emphasizes that the product of a matrix and its inverse is the identity matrix, highlighting the property of matrix inverses.  The phrase "Lots of Handwaving on details" suggests that the derivation of the inverse matrix is not shown in detail on this slide, but the result is presented. The context of the course suggests that this slide is part of a section introducing matrix operations, which are fundamental to many AI algorithms.%
\subsection{Figures}%
\begin{figure}[H]%
\centering%
\end{figure}

%
\newpage%
\section{Slide 11}%
\label{sec:Slide11}%
\subsection{Content}%

\textbf{Equations}

Systems of equations and variables can be represented as matrix operation.

For example:

$-3x + 2y = -2$
$5x - 2y = 7$

$A = \begin{pmatrix} -3 & 2 \\ 5 & -2 \end{pmatrix} \quad x = \begin{pmatrix} x \\ y \end{pmatrix} \quad b = \begin{pmatrix} -2 \\ 7 \end{pmatrix}$

$Ax = b$

$\begin{pmatrix} -3 & 2 \\ 5 & -2 \end{pmatrix} \begin{pmatrix} x \\ y \end{pmatrix} = \begin{pmatrix} -3x + 2y \\ 5x - 2y \end{pmatrix} = \begin{pmatrix} -2 \\ 7 \end{pmatrix}$

Example from https://www.cuemath.com/algebra/matrix-equation/
%
\subsection{Description}%
This slide from a course on Artificial Intelligence Basics demonstrates how systems of linear equations can be represented using matrices.  It shows a specific example of two equations with two variables, $x$ and $y$.  The equations are $-3x + 2y = -2$ and $5x - 2y = 7$.  These equations are represented in matrix form as $Ax = b$, where $A$ is a 2x2 coefficient matrix, $x$ is a 2x1 variable matrix, and $b$ is a 2x1 constant matrix.  The slide explicitly shows the matrix representations of $A$, $x$, and $b$.  The slide also shows the matrix multiplication $Ax = b$, demonstrating how the system of equations is equivalent to a matrix equation.  The example is taken from an external website (https://www.cuemath.com/algebra/matrix-equation/), indicating a reference for further learning about matrix representations of systems of equations. The context of the course suggests that this slide is part of a section introducing the mathematical tools used in AI, specifically matrix operations, which are fundamental to many AI algorithms.%
\subsection{Figures}%
\begin{figure}[H]%
\centering%
\end{figure}

%
\newpage%
\section{Slide 12}%
\label{sec:Slide12}%
\subsection{Content}%

\textbf{Solving Equations}

With invertible matrices, we can solve the system of equations as such:

$Ax = b$

$A^{-1}Ax = A^{-1}b$

Since $A^{-1}A$ is the identity matrix, $A^{-1}Ax = Ix = x = A^{-1}b$
%
\subsection{Description}%
This slide from a course on Artificial Intelligence Basics discusses solving systems of linear equations using matrix inverses.  It states that with invertible matrices, the system of equations $Ax = b$ can be solved by multiplying both sides of the equation by the inverse of matrix $A$.  The slide shows the steps: first, multiplying both sides by $A^{-1}$ to get $A^{-1}Ax = A^{-1}b$.  Then, it explains that since $A^{-1}A$ is the identity matrix ($I$), the equation simplifies to $Ix = x = A^{-1}b$.  This demonstrates how the solution vector $x$ can be found by multiplying the inverse of matrix $A$ by the constant vector $b$. The context of the course suggests that this slide is part of a section introducing matrix operations and their applications in solving systems of equations, which are fundamental to many AI algorithms.%
\subsection{Figures}%
\begin{figure}[H]%
\centering%
\end{figure}

%
\newpage%
\section{Slide 13}%
\label{sec:Slide13}%
\subsection{Content}%

\textbf{Solving the example}

$Ax = b$, in our example:

$A = \begin{pmatrix} -3 & 2 \\ 5 & -2 \end{pmatrix} \quad x = \begin{pmatrix} x \\ y \end{pmatrix} \quad b = \begin{pmatrix} -2 \\ 7 \end{pmatrix}$

$A^{-1} = \begin{pmatrix} \frac{1}{2} & \frac{1}{2} \\ \frac{5}{4} & \frac{3}{4} \end{pmatrix}$

$x = A^{-1}b = \begin{pmatrix} \frac{1}{2} & \frac{1}{2} \\ \frac{5}{4} & \frac{3}{4} \end{pmatrix} \begin{pmatrix} -2 \\ 7 \end{pmatrix} = \begin{pmatrix} \frac{1}{2}(-2) + \frac{1}{2}(7) \\ \frac{5}{4}(-2) + \frac{3}{4}(7) \end{pmatrix} = \begin{pmatrix} \frac{-2+7}{2} \\ \frac{-10+21}{4} \end{pmatrix} = \begin{pmatrix} \frac{5}{2} \\ \frac{11}{4} \end{pmatrix}$

Plugged into our original equations:

$-3(\frac{5}{2}) + 2(\frac{11}{4}) = \frac{-15+22}{4} = \frac{7}{4} \neq -2$
$5(\frac{5}{2}) - 2(\frac{11}{4}) = \frac{25-22}{4} = \frac{3}{4} \neq 7$

%
\subsection{Description}%
This slide from an Artificial Intelligence Basics course demonstrates solving a system of linear equations using matrix inversion.  The slide shows the matrix equation $Ax = b$, where $A$ is a 2x2 coefficient matrix, $x$ is a 2x1 variable matrix, and $b$ is a 2x1 constant matrix.  The given values are $A = \begin{pmatrix} -3 & 2 \\ 5 & -2 \end{pmatrix}$, $x = \begin{pmatrix} x \\ y \end{pmatrix}$, and $b = \begin{pmatrix} -2 \\ 7 \end{pmatrix}$.  The slide then calculates the inverse of matrix $A$, $A^{-1} = \begin{pmatrix} \frac{1}{2} & \frac{1}{2} \\ \frac{5}{4} & \frac{3}{4} \end{pmatrix}$.  Next, it computes the solution vector $x$ by multiplying $A^{-1}$ by $b$, resulting in $x = \begin{pmatrix} \frac{5}{2} \\ \frac{11}{4} \end{pmatrix}$.  Crucially, the slide then plugs these values back into the original equations to verify the solution.  However, the verification shows that the calculated solution does not satisfy the original equations, indicating an error in the calculation or the presentation of the slide.  The context of the course suggests that this slide is part of a section introducing matrix operations and their applications in solving systems of equations, which are fundamental to many AI algorithms.  The error in the verification highlights the importance of careful calculation and checking in mathematical problem-solving.%
\subsection{Figures}%
\begin{figure}[H]%
\centering%
\end{figure}

%
\newpage%
\section{Slide 14}%
\label{sec:Slide14}%
\subsection{Content}%

\textbf{Multiplication}

Matrix Multiplication is a bit more complicated. A matrix multiplied to a vector applies
the linear transformation from one vector to the transformed vector.
Multiplying two matrices together creates a new transformation, the equivalent of
applying one transformation after another.
The product of an $m \times n$ matrix $A$, and an $n \times p$ matrix $B$ is an $m \times p$ matrix $C$

$c_{ij} = \sum_{k=1}^{n} a_{ik} b_{kj}$

The inner dimensions must match to multiply.
%
\subsection{Description}%
This slide, part of a course on Artificial Intelligence Basics, discusses matrix multiplication. It explains that multiplying matrices is more complex than multiplying a matrix by a vector.  A matrix multiplied by a vector applies a linear transformation.  Multiplying two matrices together is equivalent to applying one transformation after another.  The slide defines the product of an $m \times n$ matrix $A$ and an $n \times p$ matrix $B$ as an $m \times p$ matrix $C$.  The formula for calculating the element $c_{ij}$ of the resulting matrix $C$ is given as $c_{ij} = \sum_{k=1}^{n} a_{ik} b_{kj}$.  The slide emphasizes that the inner dimensions (the number of columns in the first matrix and the number of rows in the second matrix) must be equal for the multiplication to be possible.  The context of the course suggests that this slide is part of a section introducing matrix operations, which are fundamental to many AI algorithms.%
\subsection{Figures}%
\begin{figure}[H]%
\centering%
\end{figure}

%
\newpage%
\section{Slide 16}%
\label{sec:Slide16}%
\subsection{Content}%

\textbf{Transposition}

A fundamental operation in linear algebra for manipulating data

Given $A = \begin{pmatrix} 1 & 2 & 3 \\ 4 & 5 & 6 \end{pmatrix}$

$A^T = \begin{pmatrix} 1 & 4 \\ 2 & 5 \\ 3 & 6 \end{pmatrix}$
%
\subsection{Description}%
This slide, part of a course on Artificial Intelligence Basics, introduces the concept of matrix transposition.  It states that transposition is a fundamental operation in linear algebra used for manipulating data.  A matrix $A$ is given as a 2x3 matrix: $A = \begin{pmatrix} 1 & 2 & 3 \\ 4 & 5 & 6 \end{pmatrix}$.  The slide then shows the transpose of this matrix, $A^T$, which is a 3x2 matrix: $A^T = \begin{pmatrix} 1 & 4 \\ 2 & 5 \\ 3 & 6 \end{pmatrix}$.  The context of the course suggests that this slide is part of a section introducing matrix operations, which are fundamental to many AI algorithms.  The example demonstrates how to switch rows and columns of a matrix to obtain its transpose.%
\subsection{Figures}%
\begin{figure}[H]%
\centering%
\end{figure}

%
\newpage%
\section{Slide 17}%
\label{sec:Slide17}%
\subsection{Content}%
%
\subsection{Description}%
This slide, part of a course on Artificial Intelligence Basics, presents a diagram of a perceptron, a fundamental type of artificial neural network.  The diagram visually depicts the structure of a perceptron.  Input values $x_1, x_2, x_3, \dots, x_n$ are connected to a summation node ($\sum$).  Each input is weighted by a corresponding weight $W_1, W_2, W_3, \dots, W_n$.  The weighted inputs are summed at the summation node.  The output of the perceptron, $out(t)$, is determined by a threshold function (represented by the $\sqcap$ symbol) applied to the sum of the weighted inputs.  The perceptron's output is influenced by a bias term, $v_0(t)$, which is set to a constant $\theta$.  The diagram clearly illustrates the flow of information from input to output, highlighting the weighted connections and the summation process.  The context of the course suggests that this slide introduces the perceptron as a basic building block in artificial neural networks, which are crucial for various AI tasks.%
\subsection{Figures}%
\begin{figure}[H]%
\centering%
\begin{subfigure}[b]{0.15\textwidth}%
\centering%
\includegraphics[width=\textwidth]{/Users/ashoksaravanan/Coding/ScribeLec/server/output/CS 243/lectures/04L1-Linear Algebra Primer/figures/17.1.png}%
\caption{}%
\end{subfigure}%
\hspace{0.05\textwidth}%
\begin{subfigure}[b]{0.15\textwidth}%
\centering%

\includegraphics[width=\textwidth]{/Users/ashoksaravanan/Coding/ScribeLec/server/output/CS 243/lectures/04L1-Linear Algebra Primer/figures/17.2.png}%
\caption{}%
\end{subfigure}%
\end{figure}

%
\newpage%
\section{Slide 19}%
\label{sec:Slide19}%
\subsection{Content}%

A \ brief \ mention \ on \ Calculus \\
There \ is \ calculus \ in \ the \ course. \\
\begin{itemize}
\item You \ should \ be \ familiar \ with \ derivatives
\item We \ will \ touch \ on \ partial \ derivatives
\item We \ will \ touch \ on \ max/min/gradiants
\end{itemize}
%
\subsection{Description}%
This slide provides a brief overview of calculus topics that will be covered in the course.  It states that calculus is part of the curriculum.  The bullet points list the specific calculus concepts that students should be familiar with or will be introduced to: derivatives, partial derivatives, and maximization/minimization using gradients.  The context of the course (Artificial Intelligence Basics) suggests that these calculus concepts will be applied to AI algorithms and models.%
\subsection{Figures}%
\begin{figure}[H]%
\centering%
\end{figure}

%
\newpage%
\section{Slide 20}%
\label{sec:Slide20}%
\subsection{Content}%
%
\subsection{Description}%
This slide displays the logistic sigmoid function and its derivative.  The slide highlights a key function, σ(x), defined as 1 / (1 + e<sup>-x</sup>).  The slide then presents the derivative of this function, dσ(x)/dx, which is calculated as e<sup>-x</sup> / (1 + e<sup>-x</sup>)<sup>2</sup>.  The URL at the bottom of the slide likely points to an external resource providing further information or derivation of this result.  The context of the course (Artificial Intelligence Basics) suggests that the logistic sigmoid function and its derivative are important components of various AI models, particularly in activation functions and classification tasks.%
\subsection{Figures}%
\begin{figure}[H]%
\centering%
\end{figure}

%
\newpage%
\section{Slide 21}%
\label{sec:Slide21}%
\subsection{Content}%

\textbf{The derivative} \\
\frac{d}{dx}\sigma(x) = \frac{d}{dx}\left(\frac{1}{1+e^{-x}}\right) = \frac{d}{dx}(1+e^{-x})^{-1} \\
= -(1+e^{-x})^{-2} \cdot \frac{d}{dx}(1+e^{-x}) \\
= -(1+e^{-x})^{-2} \cdot \left(\frac{d}{dx}(1) + \frac{d}{dx}(e^{-x})\right) \\
= -(1+e^{-x})^{-2} \cdot \left(0 + e^{-x} \cdot \frac{d}{dx}(-x)\right) \\
= -(1+e^{-x})^{-2} \cdot \left(e^{-x} \cdot (-1)\right) \\
= \frac{e^{-x}}{(1+e^{-x})^2}
%
\subsection{Description}%
This slide shows the derivation of the derivative of the logistic sigmoid function, σ(x) = 1 / (1 + e<sup>-x</sup>).  The steps are clearly laid out, starting with the application of the chain rule to the expression (1 + e<sup>-x</sup>)<sup>-1</sup>.  The derivative of 1 is 0, and the derivative of e<sup>-x</sup> is -e<sup>-x</sup>.  The final result is the derivative of the sigmoid function, which is e<sup>-x</sup> / (1 + e<sup>-x</sup>)<sup>2</sup>.  The steps are annotated with explanations like "Rewritten with negative exponents," "Apply the sum rule," and "Chain rule again," making the derivation process easy to follow.  The URL at the bottom of the slide likely points to an external resource providing further information or derivation of this result.  The context of the course (Artificial Intelligence Basics) suggests that the logistic sigmoid function and its derivative are important components of various AI models, particularly in activation functions and classification tasks.%
\subsection{Figures}%
\begin{figure}[H]%
\centering%
\end{figure}

%
\newpage%
\end{document}